% Source: https://artofproblemsolving.com/community/c3h1704985_2005_national_mathcounts_sprint_14
% Stable identity: aops-topic-1704985-q01
In an $n$-row triangle, row $k$ (from the top) contains $k$ upward-pointing unit triangles, so there are $1+2+\cdots+n=\dfrac{n(n+1)}{2}$ upward triangles. Every rod is a side of exactly one upward triangle, so the number of rods is $3\cdot\dfrac{n(n+1)}{2}$. The connectors sit at the lattice points, which form rows of $1,2,\dots,n+1$ points, so there are $\dfrac{(n+1)(n+2)}{2}$ connectors.

For $n=2$ this gives $9$ rods and $6$ connectors, as in the figure. For $n=8$: $3\cdot36=108$ rods and $\dfrac{9\cdot10}{2}=45$ connectors, so $\boxed{153}$ pieces.

(The ``Newton's little formula'' of the official solution is the method of finite differences: the totals for $n=1,2,3,4$ are $6,15,28,45$, with first differences $9,13,17$ and constant second difference $4$, so the total is a quadratic in $n$, namely $2n^2+3n+1=(2n+1)(n+1)$; for $n=8$ this is $17\cdot9=153$.)
